\subsection{Java async-profiler}\label{java-async-profiler}

\href{https://github.com/async-profiler/async-profiler}{async-profiler} je uzorkujući profajler za HotSpot JVM. Može da meri CPU vreme, alokacije, zaključavanja, wall-clock vreme i da izveze rezultat kao flame graph ili JFR snimak. Koristan je kada \texttt{perf} vidi JVM proces, ali ne daje dovoljno dobar prikaz Java metoda, JIT-ovanog koda i runtime aktivnosti.

Primer je namerno jednostavan Java program bez Maven ili Gradle projekta. Jedna nit troši CPU tražeći proste brojeve, druga pravi alokacije, a dve niti se nadmeću oko istog zaključavanja:

\codefile{ProfilerTarget.java}{java}

Prevedimo program:

\begin{minted}{bash}
$ make
\end{minted}

Pokrenemo ga u jednom terminalu:

\begin{minted}{bash}
$ java ProfilerTarget 120
pid=12345 duration=120s
Attach async-profiler from another terminal.
\end{minted}

U drugom terminalu koristimo PID koji je program ispisao. Naziv komande može biti \texttt{asprof} ili putanja do skripte iz raspakovane async-profiler arhive.

\paragraph{CPU profil}\label{async-profiler-cpu}

CPU profil je početna analiza kada nas zanima gde JVM proces aktivno troši procesorsko vreme:

\begin{minted}{bash}
$ asprof -e cpu -d 30 -f async-cpu.html 12345
\end{minted}

U flame graph-u očekujemo da se izdvoje \texttt{cpuWorker} i \texttt{isPrime}. Širina okvira predstavlja broj uzoraka, ne trajanje jednog poziva. Širok okvir znači da je metoda često bila na steku u trenucima uzorkovanja.

\paragraph{Alokacije}\label{async-profiler-alloc}

Alokacioni profil pokazuje gde nastaju objekti i koliko bajtova prolazi kroz alokator:

\begin{minted}{bash}
$ asprof -e alloc -d 30 -f async-alloc.html 12345
\end{minted}

U ovom primeru treba da se vidi \texttt{allocationWorker}. Ako CPU profil ne pokazuje ovu metodu kao dominantnu, a alokacioni profil je pokazuje, problem nije nužno CPU petlja nego pritisak na GC i memorijski podsistem.

\paragraph{Zaključavanja}\label{async-profiler-lock}

Profil zaključavanja meri vreme izgubljeno na kontenciju:

\begin{minted}{bash}
$ asprof -e lock -d 30 -f async-lock.html 12345
\end{minted}

Ovde očekujemo trag kroz \texttt{lockWorker}, jer dve niti ulaze u isti \texttt{synchronized} blok. Ovakav nalaz treba povezati sa dizajnom deljenog stanja: optimizacija tela petlje neće pomoći ako niti većinu vremena čekaju na isti monitor.

\paragraph{Wall-clock i JFR izlaz}\label{async-profiler-wall-jfr}

Wall-clock profil uzorkuje proteklo vreme, pa može pokazati i čekanja:

\begin{minted}{bash}
$ asprof -e wall -d 30 -f async-wall.html 12345
\end{minted}

Za analizu u Java Mission Control-u ili arhiviranje merenja možemo napraviti JFR datoteku:

\begin{minted}{bash}
$ asprof -e cpu -d 30 -f async-cpu.jfr 12345
\end{minted}

\paragraph{Zamke}\label{async-profiler-zamke}

\begin{itemize}
\tightlist
\item
  Profilisati reprezentativan ulaz i dovoljno dugo izvršavanje; kratki profili su često šum.
\item
  Proveriti da li se meri CPU, alokacija, lock ili wall-clock događaj. Različiti događaji odgovaraju na različita pitanja.
\item
  U kontejnerima i restriktivnim sistemima mogu biti potrebna dodatna podešavanja za \texttt{perf\_event\_paranoid}, capability-je ili pristup procesu.
\item
  Flame graph pokazuje gde je program bio tokom uzorkovanja. Ne dokazuje sam po sebi da je promena koda bezbedna.
\end{itemize}
